\documentclass{article}
\usepackage[utf8]{inputenc}
\usepackage[english]{babel}
\usepackage{amsmath}
\usepackage{graphicx}
\usepackage{hyperref}

\title{My Article on NeoLeaf}
\author{Main Author}
\date{\today}

\begin{document}

\maketitle

\begin{abstract}
This document shows NeoLeaf at work: a LaTeX project that lives in your own GitHub repository and compiles in the browser.
\end{abstract}

\section{Introduction}
NeoLeaf lets you edit your LaTeX documents right in the browser. Your files stay in your browser and in your GitHub repository, and you decide when to push them.

\section{Mathematics}
We can write equations such as:
\begin{equation}
    E = mc^2
\end{equation}

And more elaborate integrals:
\begin{equation}
    \int_{-\infty}^{+\infty} e^{-x^2} \, dx = \sqrt{\pi}
\end{equation}

\section{Collaboration}
To invite colleagues to edit this project:
\begin{enumerate}
    \item Open the sharing dialog of the project.
    \item Add your colleague's GitHub username.
    \item Your colleague opens NeoLeaf and the project appears in their list.
\end{enumerate}

\section{Conclusion}
Working with LaTeX has never been this practical and accessible.

\end{document}
